% ----- Consignes exo 53 ----- %
\begin{td-exo}[Complexité paramétrique: Kernelization pour le problème de \textsc{Vertex Cover}]\, % 53 
	\begin{definition}
		Un problème \(\Pi\) de paramètre \(k\) est de \defemph{complexité paramétrique polynomiale} (FPT) si la complexité en temps de \(\Pi\) pour une instance de taille \(n\) est bornée par \(f(k) \cdot n^{O(1)}\) pour une certaine fonction \(f\).
	\end{definition}

	Dans la suite nous allons nous intéresser au problème de \textsc{Vertex Cover} avec pour paramètre \(k\).

	Nous allons procéder à la kernelization du problème de \textsc{Vertex Cover}.
	Pour finir, nous proposerons une méthode exacte, utilisant en temps exponentiel \(k\) pour résoudre de manière exacte le problème.
	Pour cela nous allons procéder à plusieurs réductions de l'instance initiale et après utiliser un arbre borné de recherche.

	\begin{itemize}
		\item Si \(x\) est un sommet isolé alors \(x\) n'appartion à aucun vertex cover de \(G\) et on peut le supprimer de \(G\).
		\item Si \(x\) est un sommet de degré \(1\) voisin de \(y\) alors il existe une solution optimale contenant \(y\) donc
		\begin{equation*}
			VC(G, k) = VC(G \setminus \{x, y\}, k - 1) \cup \{y\}
		\end{equation*}
		\item Si \(x\) est un sommet de degré \(\geq k + 1\), alors si \(G\) possède une solution de taille \(k\) alors \(x\) doit appartenir à cette solution. Donc
		\begin{equation*}
			VC(G, k) = VC(G \setminus \{x\}, k - 1) \cup \{x\}
		\end{equation*}
	\end{itemize}

	\begin{enumerate}
		\item Justifier les trois règles précédentes.

		\item Montrer que si \(G\) possède un vertex cover de taille \(k\) alors le graphe réduit possède au plus \(k^2 + k\) sommets et au plus \(k^2\) arêtes. 
		Pour cela:
		\begin{enumerate}
			\item Montrer que le graphe réduit possède au plus \(k^2\) arêtes.
			\item Montrer que le graphe réduit possède au plus \(k^2 + k\) sommets.
		\end{enumerate}
		
		\item Montrer que si le graphe réduit possède plus de \(k^2\) arêtes ou plus de \(k^2 + k\) sommets alors \(G\) ne possède pas de solution.

		\item Soit \({\left(x_v\right)}_{v \in V}\) une solution optimale d'une instance \((G, k)\) du problème de \textsc{Vertex Cover} obtenue par la relaxation du programme linéaire en nombres entiers et soit \(V_0, V_1, V_{1/2}\) les ensembles des solutions du programme linéaire.
	\end{enumerate}
\end{td-exo}

% ----- Solutions exo 53 ----- %
\iftoggle{showsolutions}{ 
	\begin{td-sol}[]\ % 53
		\begin{enumerate}
			\item Les trois règles sont clairement cohérentes.

			\item Il y a bien au plus \(k^2\) arêtes car toutes les arêtes sont couvertes et on a un vertex cover de taille \(k\).
			Alors, chacun des \(k\) sommets du vertex cover a au plus \(k\) voisins (par notre réduction) et donc il y a bien au plus \(k \times k = k^2\) arêtes dans le graphe réduit.

			De même, on a bien au plus \(k^2 + k\) sommets (de par le degré des sommets dans le graphe).

			\item Une solution respecte forcément les critères précédents donc si on ne respecte pas les critères on a un certificat de non existence de solution.
		\end{enumerate}
	\end{td-sol}
}{}
